\section{Introduction}
\label{sec:introduction}

Consider $98 \times 14$.
A skilled human reasoner recognises this as $(100 - 2) \times 14 = 1{,}400 - 28 = 1{,}372$, a structural shortcut that sidesteps column multiplication entirely.
When the operand shifts to $73 \times 14$, no comparable decomposition presents itself and the same reasoner reverts to standard arithmetic.
The ability to detect and exploit such structural opportunities, rather than applying a uniform procedure to every problem, is the hallmark of \emph{number sense} \citep{mcintosh1992proposed, reys1991number}.

Large language models, by contrast, tend to default to verbose chain-of-thought computation regardless of whether a problem admits a shortcut \citep{wei2022chain, kojima2022large}.
This raises a basic question: do LLMs possess latent shortcut-sensitive reasoning that can be selectively activated, or do they lack this capacity entirely?
The question has been difficult to answer because existing benchmarks are not designed to test it.
GSM8K \citep{cobbe2021gsm8k} and MATH \citep{hendrycks2021math} measure whether a model reaches the correct answer, not whether it selects an efficient strategy to get there.
A model that applies long multiplication to every problem and a model that recognises algebraic shortcuts receive identical scores, provided both arrive at the right answer.
Strategy selection, the core of number sense, is invisible to accuracy-only evaluation.

We argue that measuring number sense requires a different experimental design: one that holds problem structure constant while varying shortcut availability, then tests whether prompting interventions produce differential effects across matched item variants.
This logic parallels factorial designs in cognitive psychology, where the interaction between task condition and experimental manipulation is the quantity of interest, not main effects alone.

\paragraph{Research questions.}
We address three questions:
\begin{enumerate}[leftmargin=*,itemsep=2pt]
    \item \textbf{Selective activation.} Does number-sense prompting improve accuracy on shortcut-amenable items more than on matched controls?
    \item \textbf{Cross-model variation.} Do different models show qualitatively different shortcut profiles across categories, rather than uniformly scaling with model size?
    \item \textbf{Multi-dimensional structure.} Is number sense a single axis or a collection of dissociable competencies (e.g., magnitude estimation vs.\ cancellation vs.\ structural decomposition)?
\end{enumerate}

\paragraph{The \textsc{SenseMath} benchmark.}
To answer these questions, we introduce \textsc{SenseMath}, a controlled benchmark of 1{,}600 item families spanning 8 categories and 4 digit scales ($d \in \{2, 4, 8, 16\}$).
Six \emph{core} categories target established number-sense competencies (magnitude estimation, structural shortcuts, relative distance, cancellation, compatible numbers, landmark comparison); two \emph{extension} categories probe reasoning-level and option-level shortcuts (equation reasoning and option elimination, respectively).
Each family contains three variants (strong-shortcut, weak-shortcut, and control) that share the same template and digit scale but differ in whether a clean shortcut is available.
All items use a four-option multiple-choice format with answers extracted via \boxed{} to ensure unambiguous scoring.
An overview of the benchmark design appears in Figure~\ref{fig:overview}.

\paragraph{Key findings.}
We evaluate five models (GPT-4o-mini, GPT-4.1-mini, Qwen3-30B, Qwen3-8B, and Llama-3.1-8B) under chain-of-thought (CoT), number-sense (NC), and strict prompting conditions, totalling 72{,}000 inferences.
We find that a category-agnostic NS prompt, which mentions only ``easy calculations'' and ``mathematical intuition'' without naming any specific shortcut, produces an asymmetric effect: accuracy on shortcut-amenable items is maintained or improved while accuracy on matched controls drops, most consistently at $d{=}4$.
The per-category bar charts (Figure~\ref{fig:bar_d4}) reveal that models show selective strengths and reversals across shortcut types: Qwen models lead on magnitude estimation and structural shortcuts, while GPT-4o-mini dominates cancellation, equation reasoning, and option elimination.
No single model uniformly outperforms the others, suggesting that number sense is not a monolithic capability but a collection of partially independent dimensions.

\paragraph{Contributions.}
We make four contributions:
\begin{enumerate}[leftmargin=*,itemsep=2pt]
    \item A controlled benchmark (\textsc{SenseMath}) with 1{,}600 matched families across 8 categories and 4 scales, enabling factorial analysis of shortcut exploitation under prompt manipulation.
    \item Evidence that a category-agnostic NS prompt induces a broad shortcut-seeking mode: accuracy on shortcut-amenable items is maintained or improved while control accuracy drops. The effect is statistically reliable for two of five models (GPT-4o-mini, Qwen3-8B), with the same directional pattern in two others.
    \item A three-level evaluation framework spanning Use (solve), Judge (recognise), and Generate (create), revealing crossed dissociations in shortcut competence across models.
    \item Per-category radar analysis showing that different models excel at different shortcut types, establishing number sense as a multi-dimensional construct rather than a single ability.
\end{enumerate}

\begin{figure}[t]
    \centering
    \includegraphics[width=0.95\textwidth]{figures/fig1_overview.pdf}
    \caption{\textsc{SenseMath} overview. Each item family consists of matched strong-shortcut, weak-shortcut, and control variants administered under CoT and NS prompting. The matched design isolates selective shortcut exploitation from general prompt effects.}
    \label{fig:overview}
\end{figure}
